\usepackage{booktabs}
\usepackage{amsthm}

% Quiebre automático de líneas largas en los bloques de código:
% - Highlighting: código fuente R resaltado por pandoc
% - verbatim: salidas de R sin resaltar
\usepackage{fvextra}
\DefineVerbatimEnvironment{Highlighting}{Verbatim}{breaklines,breakanywhere,commandchars=\\\{\}}
\RecustomVerbatimEnvironment{verbatim}{Verbatim}{breaklines,breakanywhere}
\makeatletter
\def\thm@space@setup{%
  \thm@preskip=8pt plus 2pt minus 4pt
  \thm@postskip=\thm@preskip
}
\makeatother

% Traducción al español de etiquetas automáticas de LaTeX 
\renewcommand{\contentsname}{Índice general}
\renewcommand{\chaptername}{Capítulo}
\renewcommand{\figurename}{Figura}
\renewcommand{\tablename}{Cuadro}
\renewcommand{\bibname}{Bibliografía}
\renewcommand{\appendixname}{Apéndice}

% Incluir todas las entradas de book.bib en la bibliografía aunque no estén citadas
\nocite{*}

% Encabezados de página: número separado del título del capítulo/sección
\usepackage{fancyhdr}
\pagestyle{fancy}
\fancyhf{}
\fancyhead[LE]{\thepage\quad\textit{\nouppercase{\leftmark}}}
\fancyhead[RO]{\textit{\nouppercase{\rightmark}}\quad\thepage}
\renewcommand{\headrulewidth}{0.4pt}
\fancypagestyle{plain}{%
  \fancyhf{}
  \renewcommand{\headrulewidth}{0pt}
}
